\documentclass[../../../include/open-logic-section]{subfiles}

\begin{document}

\olfileid{sth}{cardinals}{cardsasords}
\olsection{Kardinal minangka Ordinal}

Sajatiné, téori kardinal kita bakal kanthi terang-terangan
nganggo téori ordinal. Tegesé, kardinal bakal kita tetepake
minangka ordinal-ordinal tartamtu. Mliginé, kita menehi
definisi iki:

\begin{defn}\ollabel{defcardinalasordinal}
Yèn $A$ bisa diurutaké becik, $\card{A}$ iku ordinal
$\gamma$ paling cilik kang nyukupi $\cardeq{A}{\gamma}$. Kanggo
sembarang ordinal $\gamma$, kita kandha yèn $\gamma$ iku
\emph{kardinal} yen lan mung yen $\gamma = \card{\gamma}$.
\end{defn}

Kita lagi wae nganggo ukara “$A$ bisa diurutaké becik”.
Kaya lumrahé ing matematika, tembung modal ing kéné mung cara
ringkes kanggo nyatakake klaim eksistensial: kandha
“$A$ bisa diurutaké becik” padha karo kandha
“ana relasi kang ngurutaké $A$ kanthi becik”.

Nanging \olref{defcardinalasordinal} nduwèni alangan. Kita
péngin supaya \emph{saben} himpunan nduwèni ukuran, yaiku supaya
$\card{A}$ ana kanggo saben~$A$. Déné definisi mau diwiwiti
karo syarat: “\emph{Yèn} $A$ bisa diurutaké becik\ldots”.
Yèn ana himpunan $A$ kang ora bisa diurutaké becik, definisi
kita ora bakal nemtokaké objek~$\card{A}$.

Mula, kanggo migunakaké \olref{defcardinalasordinal}, kita perlu
jaminan yèn saben himpunan bisa diurutaké becik. Sayangé,
jaminan iki ora ana ing~$\ZF$. Dadi, yèn kita arep nganggo
\olref{defcardinalasordinal}, kita kudu nambah aksioma anyar,
kayata:
\begin{axiom}[Pengurutan Becik]
Saben himpunan bisa diurutaké becik.
\end{axiom}
Apa Aksioma Pengurutan Becik iki bisa ditampa bakal kita rembug
ing \olref[choice][]{chap}. Nanging wiwit saiki kita bakal
nganggep aksioma iku kasedhiya. Kanthi aksioma iki, cukup
gampang mbuktèkaké yèn kardinal (miturut
\olref{defcardinalasordinal}) ana lan nduwèni sipat kang trep:

\begin{lem}\ollabel{lem:CardinalsExist}
Kanggo saben himpunan $A$:
\begin{enumerate}
	\item\ollabel{cardaexists} $\card{A}$ ana lan tunggal;
	\item\ollabel{cardaapprox} $\cardeq{\card{A}}{A}$;
	\item\ollabel{cardaidem} $\card{A}$ iku kardinal, yaiku
	$\card{A} = \card{\card{A}}$;
\end{enumerate}
\end{lem}

\begin{proof}
Tetepna $A$. Miturut Pengurutan Becik, ana pengurutan becik
$\tuple{A, R}$. Miturut
\olref[ordinals][ordtype]{thmOrdinalRepresentation}, $\tuple{A,
R}$ isomorfis karo ordinal tunggal, $\beta$. Dadi
$\cardeq{A}{\beta}$. Kanthi Induksi Transfinit, ana ordinal
$\gamma$ paling cilik kang tunggal lan nyukupi
$\cardeq{A}{\gamma}$. Mula $\card{A}
= \gamma$, kang mbuktekaké \olref{cardaexists} lan \olref{cardaapprox}.
Kanggo mbuktekaké \olref{cardaidem}, wigatèkna yèn $\delta \in \gamma$
ngakibataké $\cardless{\delta}{A}$ miturut pamilihan $\gamma$,
mula uga $\cardless{\delta}{\gamma}$ amarga relasi padha cacahé
iku relasi ekivalensi
(\olref[sfr][siz][equ]{equinumerosityisequi}). Dadi $\gamma =
\card{\gamma}$.
%Dadi, kanggo kontradiksi, upamané ana ordinal $\delta \in \gamma$ kang nyukupi $\cardeq{\gamma}{\delta}$. Mula $\cardeq{\cardeq{A}{\gamma}}{\delta}$, saéngga $\cardeq{A}{\delta}$ miturut \olref[sfr][set][equ]{equinumerosityisequi}; iki mbantah pamilihan $\gamma$ minangka ordinal paling cilik kang nyukupi $\cardeq{A}{\gamma}$.
\end{proof}

Asil sabanjuré njamin Prinsip Cantor, lan uga pratelan liya.
(Wigatèkna yèn kardinal oleh urutané saka ordinal, yaiku
$\cardfont{a} < \cardfont{b}$ yen lan mung yen $\cardfont{a} \in \cardfont{b}$. Nalika
nyatakaké iki, kita bakal nganggo aksara Fraktur kanggo objek
kang wis dingertèni minangka kardinal. Iki cukup lumrah.
Pilihan liya kang umum yaiku nganggo aksara Yunani,
amarga kardinal uga ordinal, nanging milih aksara saka
tengahing abjad, umpamané: $\kappa, \lambda$.):
\begin{lem}\ollabel{lem:CardinalsBehaveRight}
Kanggo sembarang himpunan $A$ lan $B$:
\begin{align*}
	\cardeq{A}{B} &\text{ yen lan mung yen } \card{A} = \card{B}\\
	\cardle{A}{B} &\text{ yen lan mung yen } \card{A} \leq \card{B}\\
	\cardless{A}{B}&\text{ yen lan mung yen } \card{A} < \card{B}
\end{align*}
\end{lem}

\begin{proof}
Kita bakal mbuktèkaké arah saka kiwa menyang tengen kanggo
pratelan kapindho (kasus liya padha lan dadi latihan). Mula,
gatekna diagram iki:
\begin{center}
	\begin{tikzpicture}
	\node (nodea) {$A$};
	\node[right = 6em of nodea] (nodeb) {$B$};
	\node[below = 2em of nodea] (nodecarda) {$\card{A}$};
	\node[below = 2em of nodeb] (nodecardb) {$\card{B}$};
	\draw[->] (nodea)--(nodeb);
	\draw[<->] (nodea)--(nodecarda);
	\draw[<->] (nodeb)--(nodecardb);
	\draw[->, dashed] (nodecarda)--(nodecardb);
\end{tikzpicture}
\end{center}
Panah kang pucuké loro nuduhaké !!{bijection}s, kang anané
dijamin déning \olref{lem:CardinalsExist}. Kanthi nganggep
$\cardle{A}{B}$, ana !!a{injection} $A\to B$. Saiki,
nututi panah saka $\card{A}$ menyang $A$, banjur menyang
$B$, lan pungkasane menyang $\card{B}$, kita olèh
!!a{injection} $\card{A} \to \card{B}$ (panah putus-putus).
\end{proof}\noindent Kita uga bisa nganggo \olref{lem:CardinalsBehaveRight}
kanggo mbuktèkaké manèh teorema Schr\"{o}der--Bernstein. Pratelané:
yèn $\cardle{A}{B}$ lan $\cardle{B}{A}$, mula $\cardeq{A}{B}$. Iki wis
kita nyatakaké minangka \olref[sfr][siz][sb]{thm:schroder-bernstein},
nanging bukti kapisané---kanthi gawé kang ora sithik---ana ing
\olref[sfr][infinite][card-sb]{sec}.
Saiki gatekna:

\begin{proof}[Bukti manèh Schr\"oder-Bernstein]
Yèn $\cardle{A}{B}$ lan $\cardle{B}{A}$, mula $\card{A} \leq \card{B}$
lan $\card{B} \leq \card{A}$ miturut \olref{lem:CardinalsBehaveRight}. Dadi
$\card{A} = \card{B}$ lan $\cardeq{A}{B}$ miturut Trikotomi lan
\olref{lem:CardinalsBehaveRight}.
\end{proof}
\noindent
Sanadyan bukti iki prasaja banget, kanthi meneng-meneng bukti
iki ngandelaké Panggantian (kanggo njamin
\olref[ordinals][ordtype]{thmOrdinalRepresentation}) lan
Pengurutan Becik (kanggo njamin \olref{lem:CardinalsBehaveRight}).
Kosokbaliné, bukti ing \olref[sfr][infinite][card-sb]{sec} luwih
mandiri (malah bisa ditindakake ing~$\Zminus$).

\end{document}
